% Part: computability
% Chapter: computability-theory
% Section: ce-sets

\documentclass[../../../include/open-logic-section]{subfiles}

\begin{document}

\olfileid[hi]{cmp}{thy}{ces}
\olsection{संगणनीयतः परिगणनीय समुच्चय}

\begin{defn}
कोई समुच्चय \emph{संगणनीयतः परिगणनीय} है, यदि वह रिक्त है अथवा किसी
संगणनीय फलन का परिसर है।
\end{defn}

\begin{history}
संगणनीयतः परिगणनीय समुच्चयों को \emph{पुनरावर्ततः परिगणनीय} भी कहा
जाता है। यही मूल शब्दावली है; आज दोनों नाम और उनके संक्षेप ``c.e.'' तथा
``r.e.'' सामान्यतः प्रयुक्त होते हैं।
\end{history}

\begin{explain}
आपको विचार करना चाहिए कि इस परिभाषा का क्या अर्थ है और यह शब्दावली उचित
क्यों है। विचार यह है कि यदि $S$, संगणनीय फलन~$f$ का परिसर है, तो
\[
S = \{ f(0), f(1), f(2), \dots \},
\]
और इस प्रकार $f$ को~$S$ के अवयवों का ``परिगणन'' करने वाला माना जा सकता
है। ध्यान दें कि परिभाषा के अनुसार $f$~का वर्धमान फलन होना आवश्यक नहीं,
अर्थात् परिगणन का वर्धमान क्रम में होना आवश्यक नहीं। वास्तव में $f$ का
एकैकी होना भी आवश्यक नहीं; अर्थात् परिगणन $f(0)$, $f(1)$, $f(2)$,
\dots{} में~$S$ के अवयवों की पुनरावृत्तियाँ अनुमत हैं। उदाहरणतः अचर फलन
$f(x) = 0$, समुच्चय $\{ 0
\}$ का परिगणन करता है।
\end{explain}

प्रत्येक संगणनीय समुच्चय संगणनीयतः परिगणनीय है। इसे देखने के लिए मान लें
कि $S$~संगणनीय है। यदि $S$ रिक्त है, तो परिभाषा से वह संगणनीयतः
परिगणनीय है। अन्यथा $a$ को $S$ का कोई अवयव मानें। $f$ को परिभाषित करें:
\[
f(x) =
\begin{cases}
x & \text{यदि $\Char{S}(x) = 1$} \\
a & \text{अन्यथा।}
\end{cases}
\]
तब $f$ एक संगणनीय फलन है और $S$,~$f$ का परिसर है।

\end{document}
